% option1-corpus.tex — Full corpus for EMNLP short paper
% "Length Is All You Need? A Residualized Evaluation of Training-Free Retrieval Triggers"
%
% This file is the RAW CORPUS — overlong by design.
% The final 4-page submission will be carved from this.

\documentclass[11pt]{article}
\usepackage[utf8]{inputenc}
\usepackage{amsmath,amssymb}
\usepackage{booktabs}
\usepackage{natbib}

\title{Length Is All You Need? \\
A Residualized Evaluation of Training-Free Retrieval Triggers}

\author{Anonymous EMNLP Submission}
\date{}

\begin{document}
\maketitle

% ============================================================
% ABSTRACT — placeholder, will be written last
% ============================================================
\begin{abstract}
[TO BE WRITTEN — see chat for draft]
\end{abstract}


% ============================================================
% SECTION 1: MOTIVATION + RELATED WORK (FLESHED OUT)
% ============================================================
\section{Introduction and Motivation}

Dynamic retrieval-augmented generation (RAG) systems must decide
\emph{when} to pause generation and retrieve external evidence.
This ``triggering'' decision is the critical control point:
retrieve too often and the system pays unnecessary latency;
retrieve too late and hallucinations propagate through multi-step
reasoning~\citep{Su_DRAGIN, Jiang_FLARE, Trivedi_IRCoT}.

A growing family of \textbf{training-free trigger signals} has been
proposed to solve this problem. DRAGIN~\citep{Su_DRAGIN} fires when
a composite entropy--attention score exceeds a threshold.
FLARE~\citep{Jiang_FLARE} triggers when any token in a look-ahead
sentence has probability below~$\theta$.
SeaKR~\citep{Yao_SeaKR} uses the Gram determinant of $K{=}20$
stochastic hidden-state samples.
SelfCheckGPT~\citep{Manakul_SelfCheckGPT} measures $K$-sample text
consistency via NLI or $n$-gram surprisal.
Trained alternatives exist---Self-RAG~\citep{Asai_SelfRAG} learns
reflection tokens, MBA-RAG~\citep{Tang_MBARAG} uses a multi-armed
bandit, CtrlA~\citep{Liu_CtrlA} trains internal probes---but they
require per-architecture fine-tuning and are outside our scope.

All of these methods are evaluated by the same protocol: compute a
scalar trigger score per trace, rank traces, and report AUROC against
a binary correctness label. \textbf{But is the reported AUROC
measuring genuine uncertainty, or merely response length?}

\subsection{The Length Confound: From UQ to RAG Triggers}

A convergent body of 2024--2026 work establishes that response
length is a dominant confound in LLM uncertainty evaluation:

\paragraph{The formal proof.}
\cite{Santilli_RevisitingUQEvaluationLengthBias} decompose AUROC
estimation under noisy correctness labels and prove that when both the
uncertainty score and the correctness function depend on a shared
factor---such as response length---the estimated AUROC is
\emph{systematically skewed}, not merely noisy. Rankings flip across
correctness functions. Empirically, Spearman correlations between
token length and standard UQ measures are severe:
$\rho(\text{length}, \text{Perplexity}) = -0.7$,
$\rho(\text{length}, \text{Naive Entropy}) = -0.6$,
$\rho(\text{length}, \text{Semantic Entropy}) = -0.7$.
Token Length itself is a competitive UQ-method baseline on their
leaderboard.

\paragraph{The empirical finding.}
\cite{Janiak_IllusionProgressHallucinationDetection} demonstrate that
``simple heuristics based on response length can rival complex
hallucination detection techniques.'' Switching from ROUGE to
LLM-as-Judge labels collapses Perplexity performance by 45.9\% and
EigenScore by 30.4\% on NQ-Open, SQuAD, and TriviaQA. Length-based
heuristics match or surpass Semantic Entropy on multiple cells.

\paragraph{The behavioral mechanism.}
\cite{Zhang_VerbosityVeracityVCBehavior} identify \emph{Verbosity
Compensation} (VC): LLMs generate longer, wordier responses when
uncertain. VC frequency reaches 50.4\% even on GPT-4. Verbose
responses exhibit higher uncertainty across all five evaluated
datasets---including HotpotQA, MuSiQue, and
2WikiMHQA~\citep{Zhang_VerbosityVeracityVCBehavior}---with a
performance gap of 27.6\% on Qasper. This supplies the \emph{causal
story}: length is a proxy for verbosity compensation, which is a
proxy for uncertainty, which is a proxy for
incorrectness---and, we argue, for retrieval need.

\paragraph{The residualization fix.}
\cite{Vashurin_UncertaintyLineLengthInvariant} show that even
nominally length-normalized UQ measures (Perplexity, Mean Token
Entropy, Token-SAR) exhibit residual length bias, and propose
\textsc{Uncertainty-Line}: regress the uncertainty score on output
length via OLS, then use the \emph{residual} as the length-invariant
estimate:
\[
\hat{u}_{\textsc{line}}(\hat{y}, x)
= \hat{u}(\hat{y}, x) - \hat{f}(L(\hat{y})).
\]
This is post-hoc, model-agnostic, and applicable to any scalar UQ
measure. They demonstrate consistent improvements on NMT,
summarization, and mathematical reasoning.

\paragraph{The correctness-function pitfall.}
\cite{Ielanskyi_PitfallsUQEvaluationNLG} show that approximate
correctness functions (ROUGE, BLEU, BERTScore, LLM-as-judge)
disagree substantially on UQ method rankings---pairwise mutual AUROC
as low as 0.53, Spearman $\rho$ between rankings sometimes
\emph{negative} ($\rho = -0.66$ for BLEU vs.\ ROUGE on CoQA).
Length is one driver of this disagreement.

\paragraph{The acknowledged gap.}
\cite{Mao_ConfidenceOverTimeSTL} explicitly note that ``joint
generation probability suffers from severe length bias, making it
unsuitable for long-form reasoning'' but propose no residualization
fix. \cite{Chen_DynamicInstabilityLLMReasoning} cap output to 128
tokens as a length control but do not evaluate under
residualization.

\medskip
\noindent\textbf{The critical observation.} All five papers above
document the length confound in \emph{general} UQ or hallucination
detection. \textbf{None of them test whether the specific signals
used by published dynamic RAG triggers---DRAGIN, FLARE, SeaKR,
SelfCheckGPT---survive length control.}
The bridge from ``length confounds UQ evaluation'' to ``length
confounds RAG triggering'' has not been crossed.
We cross it.


\subsection{Answer-Trajectory Primitives}

Concurrent work establishes that step-to-step answer dynamics during
chain-of-thought (CoT) reasoning carry diagnostic signal:

\cite{Datta_DecideEarlyExplainLater} define \emph{forced answer
completion}: at each reasoning prefix, append a closing token and
sample an intermediate answer $A_i$, forming an answer trajectory
$T = (A_0, \ldots, A_n)$. They define
$\text{AnswerSwitches}(T) = \sum_i \mathbb{1}[A_i \ne A_{i+1}]$
and report that 68\% of Qwen3-4B traces never change answer, while
the model generates $\sim$760 tokens after the final switch.
Application: early stopping (saving compute), not retrieval
triggering.

\cite{Wang_ChainOfProbeCoTStepByStep} independently find the same
phenomenon via \emph{Chain-of-Probe}: append a fixed probing string
after each step and read the softmax over a closed answer set. Traces
that change answer during reasoning have $>$20\% lower accuracy than
early-commit traces.

\cite{Chen_DynamicInstabilityLLMReasoning} define a per-token
instability signal $I_t = \text{JSD}(\tilde{p}_t,
\tilde{p}_{t-1}) + \lambda H_t$ and show that \emph{timing} matters
more than magnitude: early peaks ($\rho < 0.25$) correlate with 46\%
accuracy (corrective instability), late peaks ($\rho > 0.5$) with
14\% accuracy (destructive instability)---a $>$3$\times$ accuracy
gap conditional on timing alone.

\paragraph{Terminological disambiguation.}
\cite{Zhang_AnswerAgreementRepresentationShaping} use ``answer
stability'' to denote robustness of a decoded answer to small
\emph{latent perturbations} at the trace boundary---a supervised,
representation-shaping approach achieving 86.6\% AUROC on TruthfulQA.
Our use of ``answer instability'' denotes step-to-step disagreement
under $K$-sample boundary-completion \emph{sampling}---a
training-free, decode-time signal. These are categorically different
operationalizations of a shared intuition.


\subsection{Stepwise Uncertainty in CoT}

The broader trend toward temporal modeling of LLM reasoning
uncertainty is now a multi-paper research front:
\cite{Lu_StreamingHallucinationDetectionCoT} train step-level +
prefix-level probes on hidden states achieving 87--93\% AUROC
(supervised upper bound on Qwen2.5-7B);
\cite{Mao_ConfidenceOverTimeSTL} mine discriminative Signal Temporal
Logic formulas over stepwise confidence, achieving AUROC $\sim$0.62
(parameter-free, but requiring formula mining);
\cite{Datta_DecideEarlyExplainLater} propose hold-for-$k$ causal
smoothing for online monitoring.
Our work sits in the \emph{parameter-free, training-free} sub-family
of this wave, with the specific application of retrieval triggering.


\subsection{Our Contributions}

We make three contributions:

\begin{enumerate}
\item \textbf{Evaluation protocol.} We adopt and extend the
\textsc{Uncertainty-Line} protocol of
\cite{Vashurin_UncertaintyLineLengthInvariant} to dynamic-RAG
triggering. Our extension: (i)~a new task family (multi-hop QA with
paired oracle re-decode), (ii)~a new target
($\texttt{lift\_em} = +1$: ``would retrieval have flipped this
answer from wrong to right?''), and (iii)~a new signal pool
(DRAGIN, FLARE, SeaKR, AID, etc.). The paired-oracle design
partially cancels correctness-function bias
\citep{Ielanskyi_PitfallsUQEvaluationNLG} because both arms share
the same EM function.

\item \textbf{Trigger-pool audit.} Under this protocol on 2--3
architectures (Phi-4-mini, Qwen2.5-7B, Gemma-4-26B-A4B; matched at
$\sim$3--4B active parameters) $\times$ 2 multi-hop QA datasets
(HotpotQA, 2WikiMHQA) at $n{=}200$/cell, we evaluate 16 trigger
signals. \textbf{Most collapse under length residualization.} Of 14
non-AID signals, 11 fall to chance-level (residualized
AUROC~$\leq 0.5$) on $\geq$2 of 4 evaluation cells.

\item \textbf{Answer Instability Detection (AID).} Adapting the
forced-answer-completion primitive of
\cite{Datta_DecideEarlyExplainLater} to retrieval triggering, we
define 4 AID aggregations over step-level boundary-completion
answer trajectories. The AID family achieves the highest
length-residualized AUROC on all 4 evaluated cells, but the
optimal aggregation is architecture-conditional.
$K{=}5$ multi-sample shows no advantage over $K{=}1$
(Datta-style), simplifying deployment.
\end{enumerate}

\noindent Our framing aligns the length-confound critique
(Santilli/Janiak/Vashurin) with the answer-trajectory primitive
(Datta/Wang/Chen) in a new application domain (RAG triggering)
where neither community has yet looked.


% ============================================================
% SECTION 2: EXPERIMENTS
% ============================================================
\section{Experiments}

\subsection{Paired-Oracle Protocol}

For each question $q$ in the evaluation set:
\begin{enumerate}
\item \textbf{Closed-book decode}: greedy generation with no
retrieval context, producing answer $a_c$ and trace $\tau_c$.
\item \textbf{Oracle-retrieval re-decode}: greedy generation with
gold supporting paragraphs injected into the prompt, producing
$a_o$ and trace $\tau_o$.
\item \textbf{Target label}:
$\texttt{lift\_em}(q) = \mathbb{1}[\text{EM}(a_o) >
\text{EM}(a_c)]$.
\end{enumerate}

This target answers the deployment-relevant question: ``would
retrieval have helped?'' The shared EM function partially cancels
correctness-function bias~\citep{Ielanskyi_PitfallsUQEvaluationNLG}.
All trigger signals are extracted from the \emph{closed-book} trace
$\tau_c$ only.

\paragraph{Connection to \S1.} This is the first application of
paired-oracle evaluation to dynamic-RAG trigger signals. Prior
trigger papers~\citep{Su_DRAGIN, Jiang_FLARE, Yao_SeaKR} evaluate
end-task EM gain from deployed retrieval, conflating trigger quality
with retriever quality. Our protocol isolates the trigger.


\subsection{Length Residualization}

For each trigger signal $s$ and cohort of traces:
\[
s_{\text{resid}} = s - (\hat{\alpha} + \hat{\beta} \cdot
\log(\texttt{tokens} + 1))
\]
where $\hat{\alpha}, \hat{\beta}$ are OLS coefficients fit on the
cohort. AUROC is then evaluated on $s_{\text{resid}}$.

This is the unsupervised variant of \textsc{Uncertainty-Line}
\citep{Vashurin_UncertaintyLineLengthInvariant}. We use
$\log(\texttt{tokens}+1)$ rather than raw length to avoid
overcorrection on the long-tailed token distribution---consistent
with the caveat in~\cite{Vashurin_UncertaintyLineLengthInvariant}
\S3.

Statistical significance: paired bootstrap ($B{=}10{,}000$) of
$\Delta(\text{signal AUROC} - \text{tokens AUROC})$ with 95\% CI
and Bonferroni correction across $k$ signals.

\paragraph{Connection to \S1.} Implements the fix that
\cite{Santilli_RevisitingUQEvaluationLengthBias} prove is necessary
and \cite{Vashurin_UncertaintyLineLengthInvariant} propose---but
on a new signal family (RAG triggers) and a new target
(\texttt{lift\_em}).


\subsection{Trigger-Signal Pool}

We evaluate 16 signals spanning 5 families:

\begin{table}[h]
\centering\small
\begin{tabular}{@{}llll@{}}
\toprule
\textbf{Signal} & \textbf{Published method} & \textbf{Family} \\
\midrule
\texttt{tokens\_generated} & (baseline) & Length \\
\texttt{wall\_clock\_s} & (baseline) & Length \\
\texttt{n\_steps} & (baseline) & Length \\
\texttt{dragin\_rind\_max} & DRAGIN~\citep{Su_DRAGIN} & Attention \\
\texttt{ent.max} & General entropy & Logit \\
\texttt{etc} ($\Delta^2 H$ max) & ETC~\citep{Li_ETC} & Logit \\
\texttt{flare.max} ($1{-}\min p$) & FLARE~\citep{Jiang_FLARE} & Logit \\
\texttt{hvr.max} & SELFDOUBT-T1 & Text \\
\texttt{drift\_cosine.max} & CRV-drift & Hidden \\
\texttt{var\_l2\_window.max} & CRV-variance & Hidden \\
\texttt{hidden\_norm.range} & CRV-norm & Hidden \\
\texttt{u\_cons.max} & $K$-sample consistency & $K$-sample \\
\texttt{seakr\_k5\_logdet} & SeaKR~\citep{Yao_SeaKR} proxy & Hidden-$K$ \\
\texttt{sc\_ngram} & SelfCheckGPT~\citep{Manakul_SelfCheckGPT} & $K$-text \\
\midrule
\texttt{AID flip\_rate} & Ours (adapted from~\citep{Datta_DecideEarlyExplainLater}) & Answer \\
\texttt{AID late\_flips} & Ours & Answer \\
\texttt{AID min\_majority} & Ours & Answer \\
\texttt{AID flip\_rate K=1} & Datta-style baseline & Answer \\
\bottomrule
\end{tabular}
\caption{Trigger-signal pool. Top 14: published/baseline signals.
Bottom 4: AID family (ours). SeaKR is a $K{=}5$ proxy of the full
$K{=}20$ method. SelfCheckGPT is the $n$-gram variant, not
headline NLI.}
\label{tab:pool}
\end{table}

\paragraph{Scope.} We evaluate \emph{training-free} triggers only.
Self-RAG~\citep{Asai_SelfRAG} (requires fine-tuned reflection
tokens), MBA-RAG~\citep{Tang_MBARAG} (requires per-arch bandit
training), CtrlA~\citep{Liu_CtrlA} (requires hidden-state probe
training), and ARS~\citep{Zhang_AnswerAgreementRepresentationShaping}
(requires contrastive representation shaping) are excluded by design.

\paragraph{Connection to \S1.} This pool spans the 4 major published
training-free RAG trigger methods (DRAGIN, FLARE, SeaKR,
SelfCheckGPT) plus 7 general UQ signal families. The claim ``most
collapse under length control'' is evaluated over this 14-signal
baseline pool.


\subsection{Answer Instability Detection (AID)}

At each reasoning-step boundary (spaCy sentence + 32-token
fallback), we sample $K$ boundary completions, extract answers via
regex ($\texttt{r"\textbackslash bAnswer\textbackslash s*:
\textbackslash s*(.+?)"}$), and define four aggregations:

\begin{itemize}
\item \texttt{flip\_rate}: fraction of consecutive-step
answer-set transitions (our analog of
AnswerSwitches~\citep{Datta_DecideEarlyExplainLater}).
\item \texttt{late\_flips}: flips in the last 50\% of steps only
(motivated by the corrective/destructive timing distinction of
\cite{Chen_DynamicInstabilityLLMReasoning}).
\item \texttt{min\_majority}: maximum within-step disagreement
($1 - \max_{\text{count}}/K$).
\item \texttt{flip\_rate K=1}: first sample only (Datta-style
single-sample baseline).
\end{itemize}

\paragraph{$K{=}5$ vs.\ $K{=}1$ ablation.} Across $n{=}410$
valid paired traces (R1 Phi + Qwen), $K{=}5$ shows no advantage
over $K{=}1$ on \texttt{lift\_em}: CIs straddle zero on both
architectures. This is consistent with the trajectory-method
literature~\citep{Datta_DecideEarlyExplainLater,
Wang_ChainOfProbeCoTStepByStep, Mao_ConfidenceOverTimeSTL} trending
toward single-sample sufficiency.
We recommend $K{=}1$ for deployment.


\subsection{Architecture Bracket}

\begin{table}[h]
\centering\small
\begin{tabular}{@{}lrrl@{}}
\toprule
\textbf{Model} & \textbf{Active params} & \textbf{Year} &
\textbf{Architecture} \\
\midrule
Phi-4-mini-instruct & 3.8B & 2026 & Dense Transformer \\
Qwen2.5-7B-Instruct & 7.6B & 2025 & Dense Transformer \\
Gemma-4-26B-A4B-it & $\sim$4B active & 2026 & MoE \\
\bottomrule
\end{tabular}
\caption{Model pool. All run under matched Option-A protocol: greedy
decoding, repetition penalty 1.05, \texttt{max\_new\_tokens=512}.}
\end{table}

\paragraph{Connection to \S1.} The 3-arch pool spans dense and MoE
architectures across 2 generations. This exceeds the literature norm
(DRAGIN: LLaMA-2-7B/13B only; SeaKR: LLaMA-2-7B only; FLARE:
GPT-3.5 only). Architecture-conditional trigger selection is a direct
empirical consequence of this breadth.


\subsection{Datasets}

\begin{itemize}
\item \textbf{HotpotQA-distractor} (in-domain): multi-hop QA
requiring bridge and comparison reasoning~\citep{Yang_HotpotQA}.
$n{=}200$ paired traces per architecture.
\item \textbf{2WikiMHQA-distractor} (OOD): multi-hop QA with
inference chains, bridge, and
comparison types. $n{=}200$ paired traces per architecture.
\end{itemize}

Both datasets are used by DRAGIN, SeaKR, MBA-RAG, and
\cite{Zhang_VerbosityVeracityVCBehavior} (as LongBench), making
our results directly comparable to the published trigger literature.


\subsection{Experiment Inventory}

\begin{table}[h]
\centering\small
\begin{tabular}{@{}llrl@{}}
\toprule
\textbf{ID} & \textbf{Experiment} & \textbf{$n$} &
\textbf{Case for inclusion} \\
\midrule
V2 & Phi-only length audit (11 signals) & 393/452 &
Establishes length dominance; \\
   &                                     &         &
replicates Janiak on RAG triggers \\
P0 & 3-arch $\times$ 2-dataset pilot ($K{=}5$) & 25/cell &
Direction-of-effect; schema validation \\
P0.6 & AID discovery (10 angles) & 150 pooled &
Surfaces AID; mechanism analysis \\
R1-Phi & Headline paired audit & 200/cell &
Closed-book + oracle, full signal pool \\
R1-Qwen & Headline paired audit & 200/cell &
Cross-arch replication \\
R1-Gemma & Headline paired audit & 200/cell &
MoE architecture (in flight) \\
A1 & $K{=}1$ vs $K{=}5$ ablation & 410 valid &
Resolves Datta-extension contribution \\
\bottomrule
\end{tabular}
\caption{Experiment inventory. V2 and P0 are completed predecessor
experiments; R1 is the headline campaign; A1 is the $K$-value
ablation computed post-hoc from R1 traces.}
\end{table}

\paragraph{Connection to \S1.}
\textbf{V2} is the empirical anchor that replicates
\cite{Janiak_IllusionProgressHallucinationDetection}'s ``length
rivals everything'' finding in the RAG-trigger family: on Phi
HotpotQA, \texttt{tokens\_generated} AUROC $= 0.749$ beats every
internal signal raw and residualized; DRAGIN-RIND (resid AUROC
$0.616$) is the only non-negative survivor.

\textbf{P0.6} is the discovery experiment that surfaced AID:
pooled \texttt{flip\_rate} AUROC$_{\text{help}} = 0.733$ with
$r(\text{tokens}) = 0.008$ (zero length correlation), beating
DRAGIN ($0.603$) by $+0.13$.

\textbf{R1} is the headline confirmation at publishable $n$:
AID \texttt{min\_majority} on Qwen pooled achieves resid AUROC
$0.648$, paired-bootstrap-significantly beating
\texttt{tokens\_generated}
($\Delta = +0.123$, CI $[+0.020, +0.224]$).

\textbf{A1} resolves the Datta-precedent question: $K{=}5 \approx
K{=}1$ empirically, so the multi-sample extension is not a
contribution---the \emph{application to retrieval triggering} is.


% ============================================================
% SECTION 3: RESULTS
% ============================================================
\section{Results}

% [TO BE FILLED from R1 Phi + Qwen + Gemma paired audit outputs]
% Hero table: cross-arch × cross-dataset survivor matrix
% Per-question-type spotlight: late_flips × 2Wiki/comparison
% K=5 vs K=1 ablation table
% Length-vs-signal Spearman rho table (appendix)


% ============================================================
% SECTION 4: CONCLUSIONS
% ============================================================
\section{Conclusions}

% [TO BE WRITTEN after Results are finalized]


% ============================================================
% BIBLIOGRAPHY (placeholder keys)
% ============================================================
% Use: \cite{lastName_KeywordFromPaper}
%
% Core length-confound lineage:
%   Santilli_RevisitingUQEvaluationLengthBias
%   Vashurin_UncertaintyLineLengthInvariant
%   Janiak_IllusionProgressHallucinationDetection
%   Zhang_VerbosityVeracityVCBehavior
%   Ielanskyi_PitfallsUQEvaluationNLG
%   Mao_ConfidenceOverTimeSTL
%
% Answer-trajectory lineage:
%   Datta_DecideEarlyExplainLater
%   Wang_ChainOfProbeCoTStepByStep
%   Chen_DynamicInstabilityLLMReasoning
%   Zhang_AnswerAgreementRepresentationShaping
%
% Streaming/temporal:
%   Lu_StreamingHallucinationDetectionCoT
%
% RAG triggers:
%   Su_DRAGIN
%   Jiang_FLARE
%   Yao_SeaKR
%   Manakul_SelfCheckGPT
%   Asai_SelfRAG
%   Tang_MBARAG
%   Trivedi_IRCoT
%   Liu_CtrlA
%   Li_ETC
%
% Datasets:
%   Yang_HotpotQA

\end{document}
